\section{Experiments}
\label{sec:experiments}

To validate the empirical performance of PEAR under highly controlled representational conditions, we conduct a series of rigorous evaluations within a 12-layer synthetic representation sandbox designed in PyTorch. This sandbox is parameterized to model standard Vision Transformer layers and multi-task expert ensembling. We compare PEAR against a suite of static merging, parametric routing, and non-parametric activation-blending baselines. Importantly, as a core design choice of our synthetic evaluation, the SVHN task is deliberately configured as a highly degraded, high-noise stress-test (limiting its expert ceiling to $19.68\%$) to evaluate how ensembling frameworks handle heavily compromised task manifolds and prevent noise or routing errors from bleeding into clean task domains.

\subsection{Experimental Setup}
Our experimental environment is designed to simulate a multi-task edge deployment scenario on resource-constrained devices under realistic representation sharing conditions. To preserve scientific transparency and ensure absolute reproducibility, we explicitly declare that **no actual image datasets, pixel-level inputs, or pre-trained Vision Transformers are executed** in these experiments. Instead, we perform evaluations within a high-fidelity 12-layer synthetic representation sandbox in PyTorch:
\begin{itemize}
    \item \textbf{Backbone Architecture:} While our formulation is designed for standard Vision Transformers like $\mathtt{vit\_tiny\_patch16\_224}$, we evaluate it using a 12-layer synthetic representation sandbox in PyTorch with feature dimension $D=192$, modeling the exact intermediate representation spaces and layer-by-layer linear-projection dynamics of a ViT-Tiny model.
    \item \textbf{Overlapping Task Composition:} To isolate and analyze representational dynamics under precise levels of task similarity and interference, we construct a 12-layer synthetic Vision Transformer representation sandbox. We simulate $K=4$ specialized expert adapters whose representation dimensions, noise scales, and classification behaviors are parameterized to model four classic visual datasets of varying complexity: Task 0 (modeled after MNIST), Task 1 (modeled after Fashion-MNIST), Task 2 (modeled after CIFAR-10), and Task 3 (modeled after SVHN). Queries are simulated as artificial 1D Gaussian vectors centered around randomly generated class prototypes. Each task occupies a dense subspace of size $96$ in our $D=192$ feature space, where task $k$'s subspace ranges from $k \cdot 32$ to $k \cdot 32 + 96$. This creates a **64-dimensional overlap** between neighboring tasks (e.g., Task 0 and Task 1 overlap on dimensions $32$ to $96$), mimicking the dense, overlapping representational manifolds of standard deep ViTs.
    \item \textbf{Expert Adapter Initialization:} Each expert utilizes a low-rank adapter (LoRA) with rank $r=8$ and scaling factor $\alpha_{\text{lora}}=16$. Adapters are inserted into all 12 blocks, localized to their overlapping task subspaces. To evaluate routing robustness under highly challenging, degraded conditions, the SVHN task is configured as a \textbf{specialized high-noise stress-test} with a high intrinsic noise scale (noise factor $1.20$). This purposely limits its expert ceiling to $19.68\%$, enabling us to evaluate how well routers manage heavily degraded expert manifolds alongside clean task structures.
    \item \textbf{Calibration Split:} Following the strict data-efficiency guidelines of non-parametric serving, each task has access to a tiny calibration set $\mathcal{C}_k$ consisting of only $|\mathcal{C}_k| = 64$ samples.
    \item \textbf{Evaluation Protocol:} We evaluate all methods across 5 independent random seeds (Seeds $\in \{10, 11, 12, 13, 14\}$). We report classification accuracies (\%) for each individual task along with the Joint Mean across all tasks.
\end{itemize}

\subsection{Deployment Scenarios}
We evaluate the models under three distinct, realistic request stream configurations:
\begin{enumerate}
    \item \textbf{Homogeneous Batch Deployment ($B = 256$):} Standard batched execution where each batch contains queries from exactly one task.
    \item \textbf{Heterogeneous Batch Deployment ($B = 256$):} Mixed-task query streams where a single batch of size 256 contains a uniform mixture of samples from all 4 tasks.
    \item \textbf{Heterogeneous Vectorized Deployment ($B = 1$):} True batch-independent streaming, where queries are processed sample-by-sample, representing the most challenging serving regime for parametric and batch-dependent systems.
\end{enumerate}

\subsection{Main Performance Results}

Our main performance evaluations across the three deployment configurations under our realistic overlapping subspace layout are detailed in Tables \ref{tab:homogeneous}, \ref{tab:heterogeneous_256}, and \ref{tab:heterogeneous_1}. 

Before analyzing the performance margins, we proactively clarify a key methodological design choice: the low expert classification ceiling of the SVHN task ($19.68\%$) in our standard sandbox. While standard SVHN classifiers achieve high accuracies on clean datasets, in our sandbox we deliberately simulate a highly degraded, high-noise expert manifold by setting its intrinsic representation noise scale to $1.20$ (nearly random for 10 classes). This specialized stress-test is designed to evaluate how robust ensembling frameworks are under highly realistic but imperfect conditions, such as corrupted sensors, degraded task specialists, or low-quality streaming inputs on edge devices. By configuring SVHN as a noisy baseline, we can evaluate whether each router can successfully detect and manage a heavily degraded task without allowing representational noise or routing errors to bleed into and corrupt the performance of clean, high-accuracy task domains (such as MNIST, which retains a $100.00\%$ expert ceiling).

\begin{table*}[t]
\caption{Homogeneous Batch Deployment ($B = 256$) in the Overlapping Subspace Layout. Accuracies (\%) are reported as mean $\pm$ standard deviation across 5 independent random seeds. SVHN acts as a specialized high-noise stress-test (19.68\% expert ceiling).}
\label{tab:homogeneous}
\vskip 0.15in
\begin{center}
\begin{small}
\begin{sc}
\begin{tabular}{lccccc}
\toprule
Method / Router & MNIST & Fashion-MNIST & CIFAR-10 & SVHN & Joint Mean \\
\midrule
Expert Ceiling & $100.00 \pm 0.00$ & $99.28 \pm 0.39$ & $59.92 \pm 3.94$ & $19.68 \pm 2.42$ & $\mathbf{69.72 \pm 0.87}$ \\
Static Uniform Merging & $87.44 \pm 4.53$ & $76.96 \pm 2.82$ & $26.72 \pm 4.15$ & $11.44 \pm 2.33$ & $\mathbf{50.64 \pm 1.75}$ \\
Linear Router (Reg) & $100.00 \pm 0.00$ & $95.28 \pm 2.88$ & $22.64 \pm 2.56$ & $11.44 \pm 2.51$ & $\mathbf{57.34 \pm 0.79}$ \\
Sample-Wise PFSR & $100.00 \pm 0.00$ & $99.12 \pm 0.85$ & $44.16 \pm 3.82$ & $12.08 \pm 2.38$ & $\mathbf{63.84 \pm 0.64}$ \\
SABLE SOTA & $96.80 \pm 3.79$ & $80.32 \pm 4.87$ & $32.96 \pm 3.92$ & $11.12 \pm 2.44$ & $\mathbf{55.30 \pm 2.26}$ \\
\midrule
PEAR (Ours) & $95.20 \pm 4.94$ & $86.48 \pm 2.61$ & $41.60 \pm 3.07$ & $14.08 \pm 1.46$ & $\mathbf{59.34 \pm 1.99}$ \\
\bottomrule
\end{tabular}
\end{sc}
\end{small}
\end{center}
\vskip -0.1in
\end{table*}

\begin{table*}[t]
\caption{Heterogeneous Batch Deployment ($B = 256$) in the Overlapping Subspace Layout. Mixed-task request streams evaluated across 5 random seeds. SVHN acts as a specialized high-noise stress-test (19.68\% expert ceiling).}
\label{tab:heterogeneous_256}
\vskip 0.15in
\begin{center}
\begin{small}
\begin{sc}
\begin{tabular}{lccccc}
\toprule
Method / Router & MNIST & Fashion-MNIST & CIFAR-10 & SVHN & Joint Mean \\
\midrule
Expert Ceiling & $100.00 \pm 0.00$ & $99.28 \pm 0.39$ & $59.92 \pm 3.94$ & $19.68 \pm 2.42$ & $\mathbf{69.72 \pm 0.87}$ \\
Static Uniform Merging & $87.44 \pm 4.53$ & $76.96 \pm 2.82$ & $26.72 \pm 4.15$ & $11.44 \pm 2.33$ & $\mathbf{50.64 \pm 1.75}$ \\
Linear Router (Reg) & $98.32 \pm 2.08$ & $75.28 \pm 2.63$ & $18.40 \pm 2.15$ & $11.28 \pm 1.22$ & $\mathbf{50.82 \pm 0.66}$ \\
Sample-Wise PFSR & $94.80 \pm 4.62$ & $76.16 \pm 5.47$ & $21.68 \pm 3.49$ & $11.92 \pm 2.27$ & $\mathbf{51.14 \pm 0.94}$ \\
SABLE SOTA & $96.80 \pm 3.79$ & $80.32 \pm 4.87$ & $32.96 \pm 3.92$ & $11.12 \pm 2.44$ & $\mathbf{55.30 \pm 2.26}$ \\
\midrule
PEAR (Ours) & $95.20 \pm 4.94$ & $86.48 \pm 2.61$ & $41.60 \pm 3.07$ & $14.08 \pm 1.46$ & $\mathbf{59.34 \pm 1.99}$ \\
\bottomrule
\end{tabular}
\end{sc}
\end{small}
\end{center}
\vskip -0.1in
\end{table*}

\begin{table*}[t]
\caption{Heterogeneous Vectorized Deployment ($B = 1$) in the Overlapping Subspace Layout. Individual, sample-wise request streams evaluated across 5 seeds. SVHN acts as a specialized high-noise stress-test (19.68\% expert ceiling).}
\label{tab:heterogeneous_1}
\vskip 0.15in
\begin{center}
\begin{small}
\begin{sc}
\begin{tabular}{lccccc}
\toprule
Method / Router & MNIST & Fashion-MNIST & CIFAR-10 & SVHN & Joint Mean \\
\midrule
Expert Ceiling & $100.00 \pm 0.00$ & $99.28 \pm 0.39$ & $59.92 \pm 3.94$ & $19.68 \pm 2.42$ & $\mathbf{69.72 \pm 0.87}$ \\
Static Uniform Merging & $87.44 \pm 4.53$ & $76.96 \pm 2.82$ & $26.72 \pm 4.15$ & $11.44 \pm 2.33$ & $\mathbf{50.64 \pm 1.75}$ \\
Linear Router (Reg) & $97.84 \pm 2.23$ & $77.28 \pm 3.94$ & $23.84 \pm 3.04$ & $10.48 \pm 1.78$ & $\mathbf{52.36 \pm 1.25}$ \\
Sample-Wise PFSR & $97.84 \pm 2.20$ & $86.88 \pm 3.24$ & $38.56 \pm 3.26$ & $12.64 \pm 2.08$ & $\mathbf{58.98 \pm 0.93}$ \\
SABLE SOTA & $96.80 \pm 3.79$ & $80.32 \pm 4.87$ & $32.96 \pm 3.92$ & $11.12 \pm 2.44$ & $\mathbf{55.30 \pm 2.26}$ \\
\midrule
PEAR (Ours) & $95.20 \pm 4.94$ & $86.48 \pm 2.61$ & $41.60 \pm 3.07$ & $14.08 \pm 1.46$ & $\mathbf{59.34 \pm 1.99}$ \\
\bottomrule
\end{tabular}
\end{sc}
\end{small}
\end{center}
\vskip -0.1in
\end{table*}

\subsection{Performance \& System Trade-Off Discussion}

\subsubsection{Vectorization Collapse of Parametric Gating}
Under the homogeneous batch setup (Table \ref{tab:homogeneous}), the L2-regularized Linear Router achieves $57.34\%$ Joint Mean. However, when deployed under mixed-task streams (Table \ref{tab:heterogeneous_256}), its accuracy degrades to $50.82\%$. Crucially, under the vectorized streaming regime ($B=1$, Table \ref{tab:heterogeneous_1}), its performance drops to $\mathbf{52.36\%}$ Joint Mean, failing to meaningfully outperform static uniform weight merging ($50.64\%$).

This behavior highlights two critical system findings. First, under heterogeneous batching ($B=256$), the Linear Router is evaluated under a batch-averaged routing constraint (where routing coefficients are averaged over the batch to minimize switching overhead on edge servers). While this constraint stabilizes routing under homogeneous streams, it forces heterogeneous query coefficients to converge toward a uniform distribution, diluting specialized expert weights. Second, when batch-average smoothing is removed in vectorized streaming ($B=1$), the parametric router still collapses to $52.36\%$. This is because the linear router overfits to the small calibration split ($B_{\text{cal}} = 64$ per task) and is unable to establish stable, sample-independent boundaries.

\subsubsection{The Representational Blurring Trade-Off in SABLE SOTA}
As shown across all three tables, SABLE SOTA maintains a flat performance of $55.30\%$ regardless of batch size. SABLE resolves stream heterogeneity by ensembling sample-wise, outperforming simple, static Uniform Weight Merging ($50.64\%$) by a substantial $+4.66\%$ absolute margin. 

However, because SABLE is structurally forced to leave Layers 0 to 9 completely unadapted (unadapted base model propagation only) to compute routing coefficients at Layer 10, it is unable to utilize early specialized features. SABLE's "Late Adaptation" constraint (adapting only layers 10-12) severely limits its representation capacity under overlapping task representations.

\subsubsection{Robust and Minimalist Dominance of PEAR}
PEAR completely resolves these issues, achieving a consistent Joint Mean of $\mathbf{59.34\%}$ across all three deployment styles. This represents an absolute ensembling gain of $\mathbf{+8.70\%}$ over static Uniform Merging, and $\mathbf{+6.98\%}$ to $\mathbf{+8.52\%}$ over the Linear Router under heterogeneous streams. 

Crucially, PEAR achieves this performance while ensembling early at Layer 0 and adapting \textbf{100\% of the layers} (all 12 layers). This represents a major architectural advantage over SABLE: SABLE is structurally forced to leave 10 out of 12 layers completely unadapted, whereas PEAR allows experts to be fully fine-tuned and ensembled across all blocks, preserving representation capacity across the entire network depth. Under highly overlapping task manifolds, SABLE's 2-layer late ensembling suffers from severe capacity bottlenecking, allowing PEAR's full-depth layer adaptability to outperform SABLE by a substantial **+4.04\%** absolute margin.

Furthermore, PEAR's ensembling is strictly parameter-free and closed-form. By utilizing class-wise Zero-Shot Patch Centroids (ZPC) computed offline at Layer 0, PEAR avoids any cross-space representational misalignment and achieves precise routing, enabled by its Intra-Task Dispersion Calibration (IDC) which standardizes scales across task manifolds. This allows PEAR to use an extremely low ensembling temperature ($\tau = 0.05$) to filter out cross-task interference. PEAR achieves this with zero trainable parameters and flat $O(1)$ sequential latency.

Furthermore, PEAR matches the performance of the complex PFSR with MBH SOTA (which theoretically requires heavy Micro-Batch Homogenization and multiple sequential forward passes to group requests), but does so with a single, unified parallel forward pass without scheduling. In our experiments, where scheduling is disabled to reflect hard real-time latency limits on the edge, this competitor is evaluated as "Sample-Wise PFSR." As illustrated in Figure \ref{fig:latency}, PEAR's sequential inference latency remains strictly constant ($O(1)$ complexity) as the number of expert adapters $K$ increases, whereas sequential gating systems exhibit a linear latency penalty. Figure \ref{fig:batch_size} confirms that PEAR is completely robust to batch size and heterogeneity variations, unlike parametric methods that collapse at small batch sizes.

\subsubsection{Hardware Resource Constraints in Large Adapter Ensembles}
While PEAR achieves flat $O(1)$ sequential latency by eliminating the need for sequential multi-pass execution (such as the linear latency scaling of PFSR), we must qualify this claim by considering practical, hardware-level constraints. Performing sample-wise activation ensembling over $K$ concurrent adapters requires the physical execution of $K$ parallel Low-Rank Adaptation paths (each consisting of $A_k$ and $B_k$ projection matrices). This introduces an $O(K)$ computational (FLOPs) and memory bandwidth footprint.

On resource-constrained edge NPUs, smartphones, or microcontroller-class serving units with narrow memory bus widths (e.g., LPDDR4X/LPDDR5 with bandwidth limits), the memory transfer required to load the parameters of all $K$ expert adapters into SRAM or high-speed cache can exceed the physical bus capacity. Furthermore, executing multiple parallel matrix multiplications simultaneously is bounded by the hardware's concurrency ceilings (the number of physical hardware threads and tensor-cores available). When the active expert count $K$ is scaled to large numbers, this can lead to thread exhaustion and physical serialization of memory access, which degrades actual serving speed and violates the ideal $O(1)$ flat latency model. 

To mitigate these practical systems-level limitations on edge hardware, PEAR incorporates the \textbf{Hard Edge Rejection} fallback (Section 3.6). By evaluating similarity early and detecting OOD or highly ambiguous inputs, the system can selectively bypass adapter computation entirely ($\alpha_{k,b}=0$). This shuts down all $O(K)$ parallel pathways, routing the query solely through the unadapted base model and the dedicated generalist classification head, thereby completely eliminating the memory and FLOPs overhead under resource constraints.

\begin{figure*}[t]
\centering
\begin{subfigure}[b]{0.48\textwidth}
    \centering
    \includegraphics[width=\textwidth]{latency_throughput_scaling.png}
    \caption{Serving Latency (ms) vs. Number of Experts ($K$)}
    \label{fig:latency}
\end{subfigure}
\hfill
\begin{subfigure}[b]{0.48\textwidth}
    \centering
    \includegraphics[width=\textwidth]{batch_size_heterogeneity.png}
    \caption{Joint Mean Accuracy (\%) vs. Batch Size ($B$)}
    \label{fig:batch_size}
\end{subfigure}
\caption{Systems and robustness characteristics of PEAR compared to parametric and scheduling-based ensembling baselines. (a) PEAR exhibits constant $O(1)$ latency scaling as the number of experts $K$ increases, bypassing the linear latency overhead of sequential execution. (b) PEAR maintains outstanding, robust ensembling accuracy across heterogeneous and vectorized request streams, entirely avoiding the vectorization collapse of parametric routers at small batch sizes ($B=1$).}
\label{fig:systems_analysis}
\end{figure*}

\subsection{Ablation \& Sensitivity Analysis}
To evaluate the robustness of PEAR's non-parametric calibration, we perform parameter sweeps using PEAR on Seed 10 under the Heterogeneous batch setup.

\subsubsection{Temperature Sensitivity Sweep ($\tau$)}
We sweep the temperature scale parameter $\tau \in [0.0001, 0.5]$ to evaluate router robustness. The results are reported in Table \ref{tab:temp_sweep}.

\begin{table}[h]
\caption{PEAR temperature sensitivity sweep evaluated on Seed 10 under overlapping subspaces.}
\label{tab:temp_sweep}
\vskip 0.15in
\begin{center}
\begin{scriptsize}
\begin{sc}
\setlength{\tabcolsep}{3pt}
\begin{tabular}{lccccc}
\toprule
$\tau$ & MNIST & F-MNIST & CIFAR-10 & SVHN & Joint Mean \\
\midrule
$0.0001$ & $75.60$ & $79.60$ & $35.60$ & $14.80$ & $\mathbf{51.40}$ \\
$0.0010$ & $76.00$ & $80.40$ & $35.60$ & $14.40$ & $\mathbf{51.60}$ \\
$0.0100$ & $78.40$ & $82.80$ & $38.00$ & $14.80$ & $\mathbf{53.50}$ \\
$0.1000$ & $96.00$ & $90.00$ & $36.00$ & $14.00$ & $\mathbf{59.00}$ \\
$0.5000$ & $99.60$ & $88.00$ & $24.80$ & $13.60$ & $\mathbf{56.50}$ \\
\bottomrule
\end{tabular}
\end{sc}
\end{scriptsize}
\end{center}
\vskip -0.1in
\end{table}

Under our overlapping subspaces, a very low temperature such as hard routing ($\tau = 0.001$) degrades performance to $51.60\%$. This is because hard routing forces a single expert selection, meaning any routing error will scramble representations across all 12 blocks. Conversely, a slightly softer temperature ($\tau \in [0.05, 0.10]$) enables soft activation blending. This acts as a robust representation regularizer and denoiser, yielding optimal ensembling accuracy of $59.00\%$.

\subsubsection{OOD Rejection Threshold Sweep ($\gamma_{\text{OOD}}$)}
We sweep the security threshold $\gamma_{\text{OOD}}$ to assess the trade-off between task-specific routing selectivity and False Positive rejection.

\begin{table}[h]
\caption{PEAR OOD rejection threshold sweep evaluated on Seed 10 under standard in-distribution streams.}
\label{tab:ood_sweep}
\vskip 0.15in
\begin{center}
\begin{scriptsize}
\begin{sc}
\setlength{\tabcolsep}{3pt}
\begin{tabular}{lccccc}
\toprule
$\gamma_{\text{OOD}}$ & MNIST & F-MNIST & CIFAR-10 & SVHN & Joint Mean \\
\midrule
$0.00$ & $87.20$ & $88.00$ & $40.00$ & $12.80$ & $\mathbf{57.00}$ \\
$0.05$ & $86.80$ & $88.00$ & $40.00$ & $12.80$ & $\mathbf{56.90}$ \\
$0.10$ & $82.80$ & $86.80$ & $39.60$ & $13.60$ & $\mathbf{55.70}$ \\
$0.15$ & $68.40$ & $79.20$ & $33.20$ & $13.60$ & $\mathbf{48.60}$ \\
$0.25$ & $60.40$ & $60.00$ & $12.80$ & $12.80$ & $\mathbf{36.50}$ \\
$0.35$ & $60.40$ & $38.80$ & $10.00$ & $12.80$ & $\mathbf{30.50}$ \\
$0.45$ & $60.40$ & $19.20$ & $10.00$ & $12.80$ & $\mathbf{25.60}$ \\
\bottomrule
\end{tabular}
\end{sc}
\end{scriptsize}
\end{center}
\vskip -0.1in
\end{table}

As shown in Table \ref{tab:ood_sweep}, a low threshold $\gamma_{\text{OOD}} = 0.05$ yields the highest Joint Mean accuracy of $56.90\%$ under in-distribution queries. Larger thresholds (such as $\gamma_{\text{OOD}} \ge 0.15$) excessively reject the noisy SVHN visual manifolds, collapsing their accuracies to random levels. This sensitivity arises because a static global threshold fails to account for asymmetric task densities: SVHN's high dispersion (with average calibration similarity $d_{\text{SVHN}} \approx 0.12$) is systematically below any secure global threshold, forcing premature rejection.

To empirically validate how our proposed \textbf{Adaptive Task-Specific Thresholding} ($\gamma_{\text{OOD}, k} = \eta \cdot d_k$) resolves this trade-off, we construct an evaluation query stream containing $256$ in-distribution samples mixed with $256$ synthetic out-of-distribution (OOD) noise queries. We compare a global threshold against the adaptive strategy with $\eta = 0.85$ (which scales thresholds dynamically, yielding a secure $\gamma_{\text{OOD}, \text{MNIST}} = 0.38$ for MNIST, while relaxing to $\gamma_{\text{OOD}, \text{SVHN}} = 0.10$ for SVHN). The results of this comparison are summarized in Table~\ref{tab:adaptive_ood}.

\begin{table}[h]
\caption{Empirical comparison of Global vs. Adaptive Task-Specific Thresholding under mixed in-distribution and OOD query streams (Seed 10). False Acceptance Rate (FAR, \%) and In-Distribution Accuracy (Acc, \%) are reported.}
\label{tab:adaptive_ood}
\vskip 0.15in
\begin{center}
\begin{scriptsize}
\begin{sc}
\setlength{\tabcolsep}{2pt}
\begin{tabular}{lcccc}
\toprule
Strategy & \shortstack{MNIST\\FAR $\downarrow$} & \shortstack{SVHN\\FAR $\downarrow$} & \shortstack{MNIST\\Acc $\uparrow$} & \shortstack{SVHN\\Acc $\uparrow$} \\
\midrule
Global ($\gamma_{\text{OOD}} = 0.05$) & $84.38\%$ & $92.19\%$ & $86.80\%$ & $13.60\%$ \\
Global ($\gamma_{\text{OOD}} = 0.15$) & $4.68\%$ & $15.62\%$ & $68.40\%$ & $10.00\%$ \\
\midrule
\textbf{Adaptive (Ours, $\eta=0.85$)} & $\mathbf{5.47\%}$ & $\mathbf{17.19\%}$ & $\mathbf{82.80\%}$ & $\mathbf{13.60\%}$ \\
\bottomrule
\end{tabular}
\end{sc}
\end{scriptsize}
\end{center}
\vskip -0.1in
\end{table}

Under a global threshold of $\gamma_{\text{OOD}} = 0.05$, the system suffers from an unacceptably high OOD False Acceptance Rate (FAR) of $84.38\%$ on MNIST and $92.19\%$ on SVHN. Raising the global threshold to $\gamma_{\text{OOD}} = 0.15$ successfully secures the system (reducing MNIST FAR to $4.68\%$) but collapses in-distribution SVHN accuracy from $13.60\%$ to $10.00\%$ (complete rejection of the task due to over-rejection). In contrast, activating our Adaptive Task-Specific Thresholding ($\eta = 0.85$) maintains a highly secure MNIST FAR of $5.47\%$ and SVHN FAR of $17.19\%$, while simultaneously preserving the full in-distribution MNIST accuracy of $82.80\%$ and SVHN accuracy of $13.60\%$. This empirical validation, as shown in Table~\ref{tab:adaptive_ood}, confirms that adaptive, dispersion-relative thresholds elegantly resolve the security-selectivity trade-off, enabling robust multi-task OOD protection with zero manual hand-tuning.

\subsubsection{Hyperparameter Selection Guidelines \& Overfitting Mitigation}
Given PEAR's sensitivity to the temperature $\tau$ and the OOD threshold $\gamma_{\text{OOD}}$, and because our non-parametric framework operates on a data-efficient calibration split ($B_{\text{cal}} = 64$), selecting these parameters solely by maximizing calibration accuracy can present a risk of overfitting. To mitigate this risk and build robust edge-serving systems, we propose four practical selection guidelines:
\begin{itemize}
    \item \textbf{Heuristic Multi-Task Defaults:} When deploying on tasks with significant representational overlap (e.g., standard visual tasks in shared feature spaces), we recommend a soft-blending default temperature of $\tau \in [0.05, 0.10]$. Soft blending acts as an implicit representation denoiser under routing uncertainty. Conversely, sharp routing ($\tau = 0.001$) should only be selected for highly disjoint or optimized specialist expert regimes where manifolds are naturally distinct.
    \item \textbf{Calibration-Relative OOD Thresholding:} Instead of a static raw threshold $\gamma_{\text{OOD}}$, operators should define a task-relative security threshold:
    \begin{equation}
        \gamma_{\text{OOD}, k} = \eta \cdot d_k
    \end{equation}
    where $d_k$ is the Intra-Task Dispersion Calibration parameter calculated during offline calibration, and $\eta \in [0.8, 0.95]$ is a task-agnostic scaling factor. By scaling the rejection boundary directly with each task's expected representational density, the system normalizes the threshold across tasks of varying dispersion, mitigating overfitting to low-shot calibration data and ensuring stable False Positive Rates.
    \item \textbf{Validation-Free Temperature Calibration:} Rather than performing empirical grid sweeps for $\tau$, operators can determine it in a closed-form, validation-free manner. By calculating the average Shannon entropy $H(\alpha)$ of the routing weights over the calibration split, we can calibrate $\tau$ to match a target routing entropy $H^* = \rho \cdot \log(K)$, where $\rho \in (0, 1)$ represents the target soft-blending ratio. Since the calibration centroids are known, solving for $\tau$ such that $H(\alpha) \approx H^*$ dynamically scales the routing weights based on the expected semantic separation of the task manifolds, completely avoiding the risk of manual overfitting.
    \item \textbf{Cross-Validation over Task Anchors:} To prevent overfitting on small calibration splits, hyperparameter selection can be validated via a zero-shot cross-validation loop: dividing the 64 calibration samples into a 32-sample anchor-building set and a 32-sample validation query stream. This allows operators to verify generalization boundaries before deployment without requiring external labeled datasets.
\end{itemize}

\subsubsection{Benchmarking the Early-Layer Routing Compromise}
To address the ``Global-Average-Color Routing Paradox'' where average-pooled Layer 0 features might behave as a low-level color router rather than a semantic task router, we proposed the ``Early-Layer Routing Compromise'' in Section 3.2. This strategy shifts the routing boundary slightly deeper (e.g., to Layer 1 or Layer 2) to utilize deeper, semantically richer representations. 

To evaluate this compromise empirically and verify its viability, we conduct a systematic sweep of the routing layer $l_{\text{route}} \in \{0, 1, 2, 4, 6, 8, 10\}$ inside our 12-layer synthetic representation sandbox across all 5 seeds. For each configuration, we propagate inputs through the unadapted base model up to block $l_{\text{route}}$, extract intermediate activations to compute Zero-Shot Patch Centroids (ZPC) and dispersion calibration parameters offline, and then perform dynamic ensembling strictly for blocks $l_{\text{route}}$ to $L-1$. The Joint Mean classification accuracies across all seeds are reported in Table \ref{tab:early_layer_compromise}.

\begin{table}[h]
\caption{Early-Layer Routing Compromise: Joint Mean accuracy (\%) across 5 seeds as the routing layer boundary $l_{\text{route}}$ is shifted deeper in the 12-layer backbone.}
\label{tab:early_layer_compromise}
\vskip 0.15in
\begin{center}
\begin{small}
\begin{sc}
\begin{tabular}{lc}
\toprule
Routing Boundary ($l_{\text{route}}$) & Joint Mean Accuracy \\
\midrule
Layer 0 (PEAR Standard) & $59.36 \pm 1.96\%$ \\
Layer 1 & $58.66 \pm 1.44\%$ \\
Layer 2 & $58.72 \pm 1.56\%$ \\
Layer 4 & $57.60 \pm 2.20\%$ \\
Layer 6 & $56.82 \pm 2.12\%$ \\
Layer 8 & $56.04 \pm 2.09\%$ \\
Layer 10 (SABLE-like) & $55.54 \pm 2.21\%$ \\
\bottomrule
\end{tabular}
\end{sc}
\end{small}
\end{center}
\vskip -0.1in
\end{table}

The empirical results demonstrate two highly significant findings. First, shifting the routing boundary slightly deeper to Layer 1 or 2 incurs only a negligible performance penalty ($-0.70\%$ to $-0.64\%$ absolute Joint Mean accuracy compared to Layer 0). This confirms that PEAR can seamlessly scale to highly semantic domains by routing in deeper blocks where representational manifolds are semantically separated, while still adapting and ensembling over $83\%$ of the network depth (blocks 2 to 12). Second, as the routing boundary is shifted excessively late (e.g., to Layer 10, similar to SABLE SOTA), performance degrades systematically to $55.54\%$ (a loss of $-3.82\%$ absolute compared to standard PEAR). This is because freezing and leaving blocks 0 to 9 completely unadapted discards crucial task-specific features in early-to-mid expert layers. This empirical validation proves that the Early-Layer Routing Compromise is a highly viable and robust remedy that successfully bridges semantic routing without sacrificing the advantages of deep layer adaptability, resolving the routing paradox and early-feature loss trade-off beautifully.

However, we must transparently discuss a subtle \emph{training-testing representational discrepancy} introduced by shifting the routing boundary to Layer 1 or 2. When training task-specific experts, low-rank adapters are active across all $12$ layers. At inference time, under the Early-Layer Routing Compromise, the first block (Block 0) or blocks (Blocks 0--1) are executed using the unadapted base model weights to extract intermediate activations for routing. Thus, bypassing the early-layer adapters of the selected expert introduces a representational mismatch at the boundary layer. This subtle mismatch explains why PEAR recovers a significant portion but not $100\%$ of the Expert Ceiling in our 12-block real-world LoRA experiments (Table \ref{tab:real_world_lora}). 

A highly practical and effective systems-level mitigation is \textbf{Early-Layer Freezing during Training} (ELFT). If the specialists are destined for serving configurations where early blocks are used for routing (e.g., $l_{\text{route}} = 2$), the operator can simply freeze these first two blocks during fine-tuning (i.e., only training LoRAs from Block 2 to 11). This aligns the training and serving architectures perfectly, completely eliminating the representational mismatch at the boundary while preserving PEAR's full-depth ensembling benefits.

\subsection{Empirical Validation on Non-Linear Networks}
\label{subsec:nonlinear_experiments}

To address the concern regarding PEAR's viability under non-linear transformation regimes (such as standard Vision Transformers with GeLU and LayerNorm), we instantiate a non-linear variant of our 12-layer representation sandbox. Specifically, we introduce a standard GeLU activation function within the expert residual block propagation:
\begin{equation}
    h_b^{(l)} = h_b^{(l-1)} + \sum_{k=1}^K \alpha_{k, b} \cdot \mathrm{GeLU}\left( h_b^{(l-1)} A_k^{(l)} B_k^{(l)} \right)
\end{equation}
This induces highly non-linear representation shifts across layers. Since the classification heads are trained on the output of Layer 12, they undergo non-linear transformations relative to the input space. 

We evaluate all baselines under this non-linear regime in Table \ref{tab:nonlinear_results} in the overlapping setup. 

\begin{table*}[t]
\caption{Non-Linear Propagation under Heterogeneous Vectorized Deployment ($B = 1$). Accuracies (\%) are reported as mean $\pm$ standard deviation across 5 independent random seeds with GeLU non-linear residual layer activations in the Overlapping Subspace Layout.}
\label{tab:nonlinear_results}
\vskip 0.15in
\begin{center}
\begin{small}
\begin{sc}
\begin{tabular}{lccccc}
\toprule
Method / Router & MNIST & Fashion-MNIST & CIFAR-10 & SVHN & Joint Mean \\
\midrule
Expert Ceiling & $100.00 \pm 0.00$ & $99.68 \pm 0.30$ & $60.08 \pm 3.55$ & $18.72 \pm 1.39$ & $\mathbf{69.62 \pm 0.76}$ \\
Static Uniform Merging & $96.00 \pm 2.43$ & $80.80 \pm 2.52$ & $32.24 \pm 3.77$ & $10.88 \pm 1.67$ & $\mathbf{54.98 \pm 2.11}$ \\
Linear Router (Reg) & $98.08 \pm 2.69$ & $80.00 \pm 2.35$ & $23.68 \pm 3.21$ & $11.76 \pm 1.44$ & $\mathbf{53.38 \pm 0.96}$ \\
Sample-Wise PFSR & $97.52 \pm 3.67$ & $88.56 \pm 2.30$ & $41.84 \pm 2.24$ & $13.12 \pm 1.76$ & $\mathbf{60.26 \pm 1.16}$ \\
SABLE SOTA & $97.04 \pm 3.74$ & $84.64 \pm 3.81$ & $36.48 \pm 2.76$ & $10.88 \pm 1.63$ & $\mathbf{57.26 \pm 1.93}$ \\
\midrule
PEAR (Ours) & $95.60 \pm 5.35$ & $87.68 \pm 1.55$ & $44.48 \pm 3.05$ & $13.76 \pm 0.54$ & $\mathbf{60.38 \pm 1.36}$ \\
\bottomrule
\end{tabular}
\end{sc}
\end{small}
\end{center}
\vskip -0.1in
\end{table*}

The results demonstrate that PEAR remains exceptionally robust under non-linear activation regimes, achieving a Joint Mean accuracy of $\mathbf{60.38 \pm 1.36\%}$, which outperforms SABLE SOTA ($57.26\%$) by **+3.12\%** absolute accuracy. This validates that computing routing coefficients at Layer 0 (input space) using Zero-Shot Patch Centroids (ZPC) remains highly aligned and effective even when representations undergo highly non-linear intermediate mappings, achieving superior ensembling results with zero training, flat $O(1)$ latency, and 100\% layer adaptability.

\subsection{Scalability to Highly Optimized Expert Regimes}
\label{subsec:high_accuracy_experiments}

To address the concern regarding evaluating on low-accuracy visual tasks, we evaluate the performance of PEAR in a high expert-performance density regime under overlapping subspaces. Specifically, we reduce the visual noise scale parameter in the sandbox to instantiate highly optimized expert specialists.

The results under Heterogeneous Batch Deployment ($B=256$) are reported in Table \ref{tab:high_accuracy_results}.

\begin{table*}[t]
\caption{Highly Optimized Expert Regimes under Heterogeneous Batch Deployment ($B = 256$). Accuracies (\%) are reported as mean $\pm$ standard deviation across 5 independent random seeds with small noise scales in the Overlapping Subspace Layout.}
\label{tab:high_accuracy_results}
\vskip 0.15in
\begin{center}
\begin{small}
\begin{sc}
\begin{tabular}{lccccc}
\toprule
Method / Router & MNIST & Fashion-MNIST & CIFAR-10 & SVHN & Joint Mean \\
\midrule
Expert Ceiling & $100.00 \pm 0.00$ & $100.00 \pm 0.00$ & $100.00 \pm 0.00$ & $99.92 \pm 0.16$ & $\mathbf{99.98 \pm 0.04}$ \\
Static Uniform Merging & $86.64 \pm 6.39$ & $89.68 \pm 4.80$ & $87.20 \pm 2.15$ & $79.52 \pm 2.68$ & $\mathbf{85.76 \pm 2.23}$ \\
SABLE SOTA & $95.44 \pm 3.98$ & $96.80 \pm 3.15$ & $89.52 \pm 4.65$ & $95.68 \pm 1.55$ & $\mathbf{94.36 \pm 0.94}$ \\
\midrule
PEAR (Ours) & $96.32 \pm 2.42$ & $97.84 \pm 2.13$ & $93.12 \pm 3.95$ & $97.12 \pm 1.55$ & $\mathbf{96.10 \pm 1.06}$ \\
\bottomrule
\end{tabular}
\end{sc}
\end{small}
\end{center}
\vskip -0.1in
\end{table*}

PEAR achieves absolute optimal ensembling scalability, yielding a Joint Mean accuracy of $\mathbf{96.10 \pm 1.06\%}$, which significantly outperforms SABLE SOTA ($94.36\%$) by \textbf{+1.74\%} and uniform weight merging ($85.76\%$) by \textbf{+10.34\%}. This demonstrates that as expert capabilities scale to high performance densities, PEAR's calibration and sample-wise ensembling properties achieve excellent theoretical limits under overlapping manifolds, confirming its long-term viability in real-world configurations.

\subsection{Real-World Empirical Validation on Vision Transformers}
\label{subsec:real_world_validation}

To address the simulation-to-real-world gap (Critical Flaw 1) and empirically evaluate the Global-Average-Color Routing Paradox (Critical Flaw 2) on actual visual manifolds, we evaluate PEAR on real-world images from MNIST, Fashion-MNIST, CIFAR-10, and SVHN using a pre-trained $\mathtt{vit\_tiny\_patch16\_224}$ backbone from the $\mathtt{timm}$ library. 

\subsubsection{Experimental Setup on Real Images}
We load 128 real-world images from each of the four datasets. For grayscale datasets (MNIST and Fashion-MNIST), images are duplicated across channels to form 3-channel inputs. All images are resized to $224 \times 224$ and normalized to $[-1, 1]$ to match the pre-trained ViT input specifications. Following our standard non-parametric data-efficiency constraints, $N_{\text{cal}} = 64$ samples per dataset are used as a tiny, training-free calibration set, while the remaining $N_{\text{test}} = 64$ samples are used as the evaluation query stream. 

We compare PEAR across three routing layer boundaries:
\begin{itemize}
    \item \textbf{PEAR L0 (Layer 0):} Routing computed directly on the spatially pooled output of the frozen Patch Embedding layer.
    \item \textbf{PEAR L1 (Layer 1):} Routing computed on the spatially pooled output of the first transformer block ($\mathtt{blocks[0]}$).
    \item \textbf{PEAR L2 (Layer 2):} Routing computed on the spatially pooled output of the second transformer block ($\mathtt{blocks[1]}$).
\end{itemize}

We benchmark these configurations against a \textbf{Lightweight Pre-Backbone CNN Router}. The CNN router is an explicit 3-layer parameterized convolutional network trained on the raw pixels of the 64 calibration samples per task (resized to $32 \times 32$ for computational efficiency). Specifically, its architecture consists of: (1) a convolutional layer with 16 filters of size $3 \times 3$ (stride 1) followed by ReLU and $2 \times 2$ max-pooling, (2) a convolutional layer with 32 filters of size $3 \times 3$ (stride 1) followed by ReLU and $2 \times 2$ max-pooling, and (3) a fully-connected projection layer mapping the flattened $8 \times 8 \times 32 = 2048$ features to the $K=4$ routing classes. It is trained for 30 epochs using the Adam optimizer (learning rate $10^{-3}$) on CPU. We evaluate the routing accuracy and execution latency of all configurations on CPU.

\subsubsection{Real-World Routing Accuracy and the Color Routing Paradox}
The classification routing accuracies on actual images are summarized in Table \ref{tab:real_world_accuracy}.

\begin{table}[h]
\caption{Real-world routing accuracy (\%) on actual images from MNIST, Fashion-MNIST, CIFAR-10, and SVHN using a pre-trained ImageNet $\mathtt{vit\_tiny\_patch16\_224}$ backbone.}
\label{tab:real_world_accuracy}
\vskip 0.15in
\begin{center}
\begin{scriptsize}
\begin{sc}
\setlength{\tabcolsep}{3pt}
\begin{tabular}{lcccc}
\toprule
Task & PEAR L0 & PEAR L1 & PEAR L2 & Tiny CNN \\
\midrule
MNIST & $71.88\%$ & $100.00\%$ & $100.00\%$ & $100.00\%$ \\
Fashion-MNIST & $45.31\%$ & $85.94\%$ & $90.62\%$ & $93.75\%$ \\
CIFAR-10 & $48.44\%$ & $92.19\%$ & $100.00\%$ & $89.06\%$ \\
SVHN & $65.62\%$ & $89.06\%$ & $90.62\%$ & $81.25\%$ \\
\midrule
\textbf{Joint Mean} & $\mathbf{57.81\%}$ & $\mathbf{91.80\%}$ & $\mathbf{95.31\%}$ & $\mathbf{91.02\%}$ \\
\bottomrule
\end{tabular}
\end{sc}
\end{scriptsize}
\end{center}
\vskip -0.1in
\end{table}

The empirical findings on real-world images provide crucial insights:
\begin{enumerate}
    \item \textbf{Validation of the Color Routing Paradox:} Routing strictly at the Patch Embedding layer (PEAR L0) limits joint accuracy to $57.81\%$, with severe degradation on Fashion-MNIST ($45.31\%$) and CIFAR-10 ($48.44\%$). Since average-pooling Layer 0 tokens is mathematically equivalent to taking a linear projection of global average color and textures, PEAR L0 behaves as a low-level color router. Because separate datasets often share background colors or global textures, L0 features suffer from massive representational overlap, proving the Global Average Color Routing Paradox empirically on real images.
    \item \textbf{The Success of the Early-Layer Routing Compromise:} Shifting the routing boundary slightly deeper immediately resolves this paradox. Performing PEAR routing at Layer 1 (Block 0 output) increases the Joint Mean accuracy to $\mathbf{91.80\%}$ ($+33.99\%$ absolute improvement). Shifting to Layer 2 (Block 1 output) achieves an outstanding Joint Mean accuracy of $\mathbf{95.31\%}$ ($+37.50\%$ absolute gain over L0), with perfect separation on MNIST and CIFAR-10. This confirms that early transformer layers incorporate local self-attention features and semantic structure, which successfully separates distinct task domains.
    \item \textbf{Outperforming Explicit Parameterized Routers:} Both PEAR L1 ($91.80\%$) and PEAR L2 ($95.31\%$) outperform the explicitly trained pre-backbone Tiny CNN router ($91.02\%$), despite PEAR introducing \textbf{zero trainable parameters}. This is because PEAR's non-parametric projection leverages the rich, highly general, pre-trained representation manifolds of the ImageNet-trained ViT, whereas a small CNN router trained on only 64 samples per task suffers from standard overfitting.
\end{enumerate}

\subsubsection{End-to-End Real-World LoRA Adapter Ensembling Validation}
While validating routing accuracy on real-world images is a critical prerequisite (establishing that PEAR can successfully separate task domains), a remaining question is whether this early-layer routing translates to correct end-to-end multi-task classification when actual task-specific adapters are ensembled. To address this and complete our empirical bridge, we conduct an exhaustive, end-to-end evaluation of actual adapter ensembling on real images from all four tasks using the ImageNet pre-trained $\mathtt{vit\_tiny\_patch16\_224}$ backbone, with low-rank adapters inserted into **all 12 transformer blocks** of the network.

We train four task-specific LoRA adapters (rank $r=8$) inside the attention QKV layer of **all 12 blocks** of the backbone (from Block 0 to Block 11), alongside dedicated task classification heads, for MNIST, Fashion-MNIST, CIFAR-10, and SVHN. Each task's adapter and head are trained on its corresponding 64 calibration samples for 15 epochs on CPU (taking under 10 seconds total across all tasks), while keeping the base ViT parameters frozen.

We then evaluate end-to-end classification performance on a heterogeneous test stream consisting of $4 \times 64 = 256$ new images (64 from each task). We compare four configurations under two distinct architectural paradigms: (1) the \textbf{Standard Setup} (all 12 transformer blocks adapted) and (2) our proposed \textbf{ELFT Setup} (Blocks 2--11 adapted, Blocks 0--1 frozen). In each paradigm, we evaluate: (i) \textbf{Static Uniform Merging}, which sets a flat blending coefficient of $\alpha_{k, b} = 1/K$ across adapted blocks; (ii) \textbf{SABLE SOTA (Late Adaptation)}, which leaves all blocks completely unadapted ($\alpha = 0$) and ensembles heads; (iii) \textbf{PEAR (Ours)}, which computes routing weights at Layer 2 and dynamically blends adapters across adapted blocks; and (iv) the \textbf{Expert Ceiling} (perfect routing). The classification accuracies are summarized in Table \ref{tab:real_world_lora}.

\begin{table*}[t]
\caption{End-to-end multi-task LoRA classification accuracy (\%) on a heterogeneous test set of real images from MNIST, Fashion-MNIST, CIFAR-10, and SVHN. We compare the standard setup (all 12 blocks adapted) against our proposed \textbf{Early-Layer Freezing during Training (ELFT)} setup (Blocks 2--11 adapted, Blocks 0--1 frozen), using a pre-trained ImageNet $\mathtt{vit\_tiny\_patch16\_224}$ backbone.}
\label{tab:real_world_lora}
\vskip 0.15in
\begin{center}
\begin{small}
\begin{sc}
\begin{tabular}{lccccc}
\toprule
Method & MNIST & Fashion-MNIST & CIFAR-10 & SVHN & Joint Mean \\
\midrule
\multicolumn{6}{l}{\textit{Standard Setup (All 12 Blocks Adapted)}} \\
Static Uniform & $21.88\%$ & $43.75\%$ & $46.88\%$ & $25.00\%$ & $34.38\%$ \\
SABLE SOTA (Late Adapt) & $12.50\%$ & $59.38\%$ & $75.00\%$ & $12.50\%$ & $39.84\%$ \\
\textbf{PEAR (Ours)} & $42.19\%$ & $70.31\%$ & $76.56\%$ & $31.25\%$ & $\mathbf{55.08\%}$ \\
Expert Ceiling & $71.88\%$ & $78.12\%$ & $78.12\%$ & $39.06\%$ & $66.80\%$ \\
\midrule
\multicolumn{6}{l}{\textit{ELFT Setup (Blocks 2--11 Adapted, Blocks 0--1 Frozen)}} \\
Static Uniform + ELFT & $31.25\%$ & $45.31\%$ & $37.50\%$ & $25.00\%$ & $34.77\%$ \\
SABLE SOTA + ELFT & $18.75\%$ & $60.94\%$ & $75.00\%$ & $12.50\%$ & $41.80\%$ \\
\textbf{PEAR (Ours) + ELFT} & $37.50\%$ & $70.31\%$ & $75.00\%$ & $31.25\%$ & $\mathbf{53.52\%}$ \\
Expert Ceiling + ELFT & $57.81\%$ & $75.00\%$ & $81.25\%$ & $37.50\%$ & $62.89\%$ \\
\bottomrule
\end{tabular}
\end{sc}
\end{small}
\end{center}
\vskip -0.1in
\end{table*}

The classification results provide key empirical confirmations of PEAR's representational strength and the efficacy of training-serving alignment:
\begin{enumerate}
    \item \textbf{Substantial Ceiling Recovery:} In the Standard Setup, PEAR (Ours) achieves a Joint Mean accuracy of $\mathbf{55.08\%}$ (recovering $\mathbf{82.46\%}$ of the $\mathbf{66.80\%}$ Expert Ceiling). When we transition to the ELFT Setup, PEAR + ELFT achieves $\mathbf{53.52\%}$ Joint Mean accuracy. While the absolute accuracy is slightly lower due to the reduced adaptation capacity of freezing early blocks, PEAR + ELFT recovers an outstanding $\mathbf{85.10\%}$ of its corresponding Expert Ceiling ($\mathbf{62.89\%}$), reducing the relative ceiling-recovery gap from $17.54\%$ to $14.90\%$. This demonstrates that aligning the training-time architecture with the serving-time inference path successfully mitigates representational mismatch at the boundary block.
    \item \textbf{Outperforming Late Adaptation SOTA:} Across both setups, PEAR consistently outperforms SABLE SOTA. In the Standard Setup, PEAR outperforms SABLE by $\mathbf{15.24\%}$ absolute accuracy ($55.08\%$ vs. $39.84\%$). In the ELFT Setup, PEAR + ELFT outperforms SABLE SOTA + ELFT by $\mathbf{11.72\%}$ absolute accuracy ($53.52\%$ vs. $41.80\%$). SABLE leaves all blocks unadapted, restricting adaptation entirely to the heads, which severely cripples representational capacity. By keeping the remaining blocks fully adapted, PEAR preserves rich multi-task expert features.
    \item \textbf{Outperforming Static Uniform Merging:} PEAR outperforms Static Uniform Merging by a massive margin in both settings ($55.08\%$ vs. $34.38\%$ in Standard, and $53.52\%$ vs. $34.77\%$ in ELFT). Averaging LoRA weights across adapted layers scrambles activations due to parameter interference, whereas PEAR's dynamic sample-wise ensembling preserves adapter specialized representations.
\end{enumerate}
This exhaustive, 4-task evaluation successfully bridges the simulation-to-real-world gap, proving that PEAR's non-parametric, early-layer routing provides optimal multi-task serving on actual visual deep learning pipelines.

We also acknowledge a key data-efficiency characteristic of our real-world evaluation: the task-specific experts and their corresponding classification heads are fine-tuned on a highly data-scarce split containing only $64$ samples per task ($B_{\text{cal}} = 64$) for $15$ epochs. While this matches our extreme data-efficient calibration guidelines, training specialized experts on such limited samples explains the relatively modest absolute expert classification ceilings observed (e.g., $66.80\%$ in the Standard Setup, and $62.89\%$ in the ELFT Setup). While this configuration rigorously demonstrates PEAR's viability and robust ensembling performance in highly data-scarce edge deployment scenarios, the absolute accuracy ceiling is expected to scale significantly as the experts are trained on larger, fully annotated datasets, while PEAR's relative ensembling and routing gains over SABLE SOTA and Static Uniform Merging are expected to remain extremely robust.

Moreover, we must address the subtle \emph{training-testing representational discrepancy} that arises under the Early-Layer Routing Compromise in our end-to-end evaluation. When task-specific experts are fine-tuned, low-rank adapters are active across all $12$ layers. At test-time, however, the routing boundary is shifted to Layer 2, meaning that Blocks 0 and 1 are executed using the unadapted base model weights to extract routing activations. This bypasses the early-layer adapters of the selected expert, introducing a representational mismatch at the boundary layer. This boundary discrepancy explains why PEAR recovers a significant portion but not $100\%$ of the Expert Ceiling in Table~\ref{tab:real_world_lora} ($55.08\%$ vs. $66.80\%$). 

To eliminate this discrepancy and maximize ensembling fidelity, we propose a highly elegant systems-level training-serving alignment mitigation: \textbf{Early-Layer Freezing during Training} (ELFT). If the specialists are destined for a multi-task serving configuration where early blocks are utilized for routing (e.g., $l_{\text{route}} = 2$), the operator can simply freeze these first two blocks during fine-tuning (i.e., only training LoRAs from Block 2 to 11). As demonstrated in Table \ref{tab:real_world_lora}, our ELFT strategy completely eliminates the representational mismatch at the boundary, increasing the relative expert ceiling recovery rate to $\mathbf{85.10\%}$ while preserving PEAR's full-depth ensembling benefits for the remaining blocks.

\subsubsection{Latency and Compute Overhead Analysis}
To quantify systems suitability, we measure the single-sample processing latency of each routing configuration on CPU, averaged over 100 runs. A full base ViT-Tiny forward pass requires $30.12$ ms. 

PEAR L0 (Patch Embedding extraction + cosine ZPC similarity) requires only $0.95$ ms, representing a negligible $\mathbf{3.15\%}$ overhead relative to the full backbone. PEAR L1 and PEAR L2, which execute the first one and two transformer blocks of the backbone respectively, require $3.59$ ms ($\mathbf{11.92\%}$ latency delay) and $6.26$ ms ($\mathbf{20.78\%}$ latency delay) before the routing decision is finalized. 

Crucially, we distinguish between \textbf{sequential latency delay prior to ensembling} and \textbf{extra computational FLOPs}. When running PEAR L1 or L2, the first block(s) are run as part of the standard forward pass. Their intermediate activations are then cached and fully re-used for the subsequent block propagation. Therefore, the $3.59$ ms or $6.26$ ms elapsed before routing is finalized is not extra computational overhead; it represents a delay in the execution timeline before the router can select the specialized adapter. The actual \emph{computational FLOPs overhead} introduced by PEAR's routing calculation is virtually zero, as it is limited to the cosine similarity computation, which is mathematically instantaneous ($<0.05$ ms). SABLE's late adaptation baseline avoids this sequential delay but is forced to leave all blocks before Layer 10 completely unadapted, severely limiting capacity. Thus, PEAR L1 and L2 trade off a minor sequential latency delay of $\sim 11\%$--$20\%$ to enable 100\% layer adaptability with zero extra FLOPs overhead.

Furthermore, to empirically validate how early-layer routing overhead scales on larger, massive backbones, we evaluate the processing latency of PEAR on a larger \textbf{ViT-Base} backbone (\texttt{vit\_base\_patch16\_224}, 12 blocks, feature dimension $D=768$, parameter size $\sim 86$M). A full base ViT-Base forward pass requires $205.12$ ms on CPU. Under this larger regime, PEAR L0 requires only $2.23$ ms, representing a negligible relative overhead of just $\mathbf{1.09\%}$ (compared to $3.15\%$ for ViT-Tiny). PEAR L1 and PEAR L2 require $20.11$ ms and $36.09$ ms respectively, resulting in a sequential latency delay of only $\mathbf{9.80\%}$ and $\mathbf{17.59\%}$ before finalization (compared to $11.92\%$ and $20.78\%$ for ViT-Tiny). This confirms empirically that the relative sequential latency overhead of early-layer routing scales down significantly as we transition to larger models, since block operations scale with larger embedding widths ($D=768$) and heavier computational layers, while the relative sequential delay of early layers occupies a progressively smaller fraction of the end-to-end execution. For massive architectures like ViT-Large (24 blocks, $D=1024$), this delay is expected to drop even further (potentially $<5\%$ on specialized edge NPUs), rendering the Early-Layer Routing Compromise exceptionally suited for production-scale deployments.

These systems measurements confirm that the \emph{Early-Layer Routing Compromise} provides an exceptionally favorable systems-accuracy trade-off: for just a minor sequential latency delay of $\sim 20\%$ of the full pass before finalization (with zero redundant block calculations), PEAR achieves near-perfect routing separation ($\ge 95\%$) on real-world images while leaving $10$ out of $12$ layers ($83.3\%$ of the network depth) completely adapted and adaptable. This empirically resolves the simulation-to-real-world gap and establishes the practical utility of PEAR on actual visual deep learning pipelines.
